% Generated by tools/edn_latex.py from reports/edn.md.
% Do not edit: change reports/edn.md and run `make edn`.
\ednentry{
  id = {EDN-80},
  title = {The project is the used-car price component, the most feasible-sounding of four ideas, chosen without a detailed comparison},
  date = {2026-09-21 at the latest (the first commit's README names car price estimation); the day itself is not recorded},
  milestone = {M1: Project Inception},
  topic = {Project Selection},
  participants = {The team, as @lukas2510 recalls it; who took part in the discussion was not recorded},
  decision = {We build a used-car price component: a model behind an API that estimates the market price of a car and the price range a buyer should expect to pay.},
  rationale = {\emph{Alternatives considered:}
\begin{itemize}[leftmargin=*, topsep=2pt]
\item \textbf{Option A (chosen): the used-car price component.}
Pros, as seen at the time: it sounded the most feasible of the four ideas, and the team found it interesting.
\newline Cons: none was weighed at the time; two showed later, the dataset's skew towards premium brands (EDN-01) and SC-04's miss on cars over 20 years (EDN-62).
\item \textbf{Option B: a route-safety model.}
Not investigated in depth: we looked for no dataset and defined no target for it, so its pros and cons were never established.
\item \textbf{Option C: a recommender.}
Not investigated in depth, as option B.
\item \textbf{Option D: a grocery idea.}
Not investigated in depth, as option B.
\end{itemize}
\par
\emph{Why:} This is @lukas2510's recollection, written down on 2026-10-12, and not a record from the time.
Four ideas were on the table, and the team did not weigh them at length.
The used-car price idea sounded the most feasible and interested us, so we checked that it was feasible and started.
The other three were not rejected for a reason of their own: they were not picked.
\par
What confirmed the choice afterwards, added on 2026-10-12 and not reasons the team had at the time:
a public dataset of 118,382 listings with a DOI existed (EDN-01; its two licence readings are in EDN-21);
price prediction is a regression with a measurable target, so the success criteria could be written as numbers (SC-01 to SC-06, EDN-06 and EDN-15);
and each LightGBM variant fits in under a minute on a laptop (EDN-69), which leaves the time for the MLOps practices the course grades.},
  ai = {Information seeking},
  response = {Used as input for further analysis},
  assessment = {As far as @lukas2510 recalls, AI did not take part in choosing among the four ideas.
It was used for the feasibility check, and its output was not taken at its word: the plan an AI assistant had drafted rested on assumptions about the data that a profile of the real files contradicted, and the project brief records the corrected picture (\href{https://github.com/mlops-2627q1-mds-upc/MLOPS_Recommenditos/pull/13}{PR \#13}).
The three reasons added afterwards are AI's, checked against the repository on 2026-10-12; a fourth, a simple API to serve, was dropped because the API does not exist yet.},
  aievidence = {Claude Code sessions of 2026-09-22 (the drafted plan checked against the data) and 2026-10-12 (@lukas2510's account).},
  evidence = {\href{https://github.com/mlops-2627q1-mds-upc/MLOPS_Recommenditos/issues/95}{issue \#95}; \href{https://github.com/mlops-2627q1-mds-upc/MLOPS_Recommenditos/blob/main/docs/docs/problem-spec.md}{problem specification}, success criteria; \href{https://github.com/mlops-2627q1-mds-upc/MLOPS_Recommenditos/blob/main/reports/figures/energy-vs-error.csv}{\texttt{reports/\allowbreak{}figures/\allowbreak{}energy-\allowbreak{}vs-\allowbreak{}error.\allowbreak{}csv}}, \texttt{fit\_\allowbreak{}seconds}.},
}
